% figures/frontier.tex -- Fig. 1 of paper/main.tex (phase L2, 2026-11-10).
% Drawn by pgfplots from paper/generated/fig_frontier_quant_<model>.dat and
% fig_frontier_sparse_<model>.dat, which scripts/make_tables.py writes (columns
% bits bits_tcr bits_mask ppl kl). No number is typed here: points, axis
% ranges and the number of points all come from those files. A panel whose two
% files have no rows yet is drawn empty with the word "pending".
% Needs \ifEsix (defined in main.tex): with E6, a fourth panel (SmolLM2-1.7B).
\ifEsix
  \def\frontierw{2.05cm}\def\frontiersep{0.75cm}\def\frontiergroup{4 by 1}
\else
  \def\frontierw{2.75cm}\def\frontiersep{0.95cm}\def\frontiergroup{3 by 1}
\fi
\pgfplotsset{
  frontier/.style={
    scale only axis, width=\frontierw, height=1.55cm,
    ymode=log, log ticks with fixed point,
    ytick={10,15,20,30,50,100,200,500,1000},
    xmin=1.3, xmax=4.25, xtick={2,3,4}, minor xtick={1.5,2.5,3.5},
    tick label style={font=\scriptsize}, label style={font=\scriptsize},
    title style={font=\scriptsize, yshift=-1.6ex},
    xlabel={bits per weight}, xlabel style={yshift=1.2ex},
    unbounded coords=jump, clip marker paths=true,
    legend style={font=\scriptsize, draw=none, column sep=5pt},
  },
}
% \frontiercheck{<file stem>}{<style name>}: style for an empty panel.
\newcommand{\frontiercheck}[2]{%
  \pgfplotstablegetrowsof{generated/fig_frontier_quant_#1.dat}\edef\frontiernq{\pgfplotsretval}%
  \pgfplotstablegetrowsof{generated/fig_frontier_sparse_#1.dat}\edef\frontierns{\pgfplotsretval}%
  \ifnum\numexpr\frontiernq+\frontierns\relax=0
    \pgfplotsset{#2/.style={ymin=10, ymax=100,
      extra description/.code={\node[font=\scriptsize\itshape, text=gray] at (0.5,0.5) {pending};}}}%
  \else
    \pgfplotsset{#2/.style={}}%
  \fi}
% \frontierplots{<file stem>}: plain quantization (black line), pruning +
% quantization (best order per cell) placed by TCr (gray circles) and with the
% 1-bit mask (blue dots), with an arrow from the first to the second.
\newcommand{\frontierplots}[1]{%
  \addplot[black, mark=square*, mark size=1.1pt, line width=0.6pt,
    restrict x to domain=2.5:4.5]% plain 2-bit is far off scale (see caption)
    table[x=bits_mask, y=ppl] {generated/fig_frontier_quant_#1.dat};
  \addplot[only marks, mark=o, gray, mark size=1.3pt]
    table[x=bits_tcr, y=ppl] {generated/fig_frontier_sparse_#1.dat};
  \addplot[blue!75!black, only marks, mark=*, mark size=1.3pt]
    table[x=bits_mask, y=ppl] {generated/fig_frontier_sparse_#1.dat};
  \addplot[forget plot, draw=none, mark=none,
    error bars/.cd, x dir=plus, x explicit, error mark=none,
    error bar style={gray, line width=0.4pt, -{Stealth[length=2.2pt]}, shorten >=1.4pt}]
    table[x=bits_tcr, y=ppl, x error expr=\thisrow{bits_mask}-\thisrow{bits_tcr}]
    {generated/fig_frontier_sparse_#1.dat};
}
\frontiercheck{SmolLM2-135M}{frontierSmolSmall}%
\frontiercheck{SmolLM2-360M}{frontierSmolMedium}%
\frontiercheck{Qwen2.5-0.5B}{frontierQwen}%
\frontiercheck{SmolLM2-1.7B}{frontierSmolLarge}%
\begin{tikzpicture}
\begin{groupplot}[frontier,
    group style={group size=\frontiergroup, horizontal sep=\frontiersep}]
\nextgroupplot[title={SmolLM2-135M}, frontierSmolSmall,
    ylabel={WikiText-2 perplexity}, ylabel style={yshift=-0.6ex}]
  \frontierplots{SmolLM2-135M}
\nextgroupplot[title={SmolLM2-360M}, frontierSmolMedium]
  \frontierplots{SmolLM2-360M}
\nextgroupplot[title={Qwen2.5-0.5B}, frontierQwen]
  \frontierplots{Qwen2.5-0.5B}
\ifEsix
\nextgroupplot[title={SmolLM2-1.7B}, frontierSmolLarge]
  \frontierplots{SmolLM2-1.7B}
\fi
\end{groupplot}
\end{tikzpicture}
